\documentclass[12pt,twoside]{article}
%\usepackage[table]{xcolor} % important to avoid options clash.
%\input{02465shared_preamble}
%\usepackage{cleveref}
\usepackage{url}
\usepackage{graphics}
\usepackage{multicol}
\usepackage{rotate}
\usepackage{rotating}
\usepackage{booktabs}
\usepackage{hyperref}
\usepackage{pifont}
\usepackage{latexsym}
\usepackage[english]{babel}
\usepackage{epstopdf}
\usepackage{etoolbox}
\usepackage{amsmath}
\usepackage{amssymb}
\usepackage{multirow,epstopdf}
\usepackage{fancyhdr}
\usepackage{booktabs}
\usepackage{xcolor}
\newcommand\redt[1]{ {\textcolor[rgb]{0.60, 0.00, 0.00}{\textbf{ #1} } } }


\newcommand{\m}[1]{\boldsymbol{ #1}}
\newcommand{\EE}{\mathbb{E}}
\DeclareMathOperator*{\argmax}{arg\,max}
\DeclareMathOperator*{\argmin}{arg\,min}
\newcommand{\yoursolution}{ \redt{(your solution here) } } 



\title{ Report 3 hand-in }
\date{ \today }
\author{Fabian (\texttt{s245632})\and  Valdemar (\texttt{s246575})}

\begin{document}
\maketitle

\begin{table}[ht!]
\caption{Attribution table. Feel free to add/remove rows and columns}
\begin{tabular}{lll}
\toprule
                                                                           & Fabian   &  Valdemar \\
\midrule
 1: Optimal policy                                                         & 60\%  & 40\%  \\
 2: Simulating a finite approximation of the optimal action-value function & 40\%  & 60\%  \\
 3: Analytically computing the optimal action-value function               & 60\%  & 40\%  \\
 4: Extend solution to all states and actions                              & 40\%  & 60\%  \\
 5: Modified reward function                                               & 60\%  & 40\%  \\
 6: Long bridge                                                            & 40\%  & 60\%  \\
\bottomrule
\end{tabular}
\end{table}

%\paragraph{Statement about collaboration:}
%Please edit this section to reflect how you have used external resources. The following statement will in most cases suffice:
%\emph{The code in the irls/project1 directory is entirely}

%\paragraph{Main report:}
Headings have been inserted in the document for readability. You only have to edit the part which says \yoursolution.

\newpage
\section{Jar-Jar at the battle of Naboo (\texttt{jarjar.py})}
\subsubsection*{{\color{red}Problem 3:  Analytically computing the optimal action-value function}}

		Using $r(s) = -|s|$ and noting that the optimal policy alternates between $s = 0$ and $s = 1$, the Bellman equations can be written as
        \begin{align*}
        V^*(0) &= \gamma V^*(1), \\
        V^*(1) &= -1 + \gamma V^*(0).
        \end{align*}
        Solving this system yields
        \begin{align*}
        V^*(0) &= \frac{-\gamma}{1 - \gamma^2}, \\
        V^*(1) &= \frac{-1}{1 - \gamma^2}.
        \end{align*}

        Applying
        \[
        Q^*(s,a) = r(s) + \gamma V^*(s + a),
        \]
        we obtain
        \begin{align*}
        Q^*(0,1) &= \frac{-\gamma}{1 - \gamma^2}, \tag{1} \\
        Q^*(1,-1) &= \frac{-1}{1 - \gamma^2}. \tag{2}
        \end{align*}

        Therefore,
        \[
        Q^*(1,-1) \leq Q^*(0,1) \leq 0 \quad \text{for all } 0 \leq \gamma < 1,
        \]
        confirming that being closer to the origin is preferable and that the agent always accumulates negative reward.


\subsubsection*{{\color{red}Problem 5:  Modified reward function}}

Using the fact that the expected reward $p \cdot r_0 + (1-p) \cdot 2r_0 = (2-p)r_0$ is the same for all states and actions, and solving the Bellman equation
  $v_\pi(s) = (2-p)r_0 + \gamma v_\pi(s)$, we get that $v_\pi(s)$ is constant across all states and policies:

  $$v_\pi(s) = \frac{(2-p),r_0}{1-\gamma}.$$


\section{Optimal escape}
\subsubsection*{{\color{red}Problem 6:  Long bridge}}
Let $c>0$ denote the step cost, so the reward for each ordinary step is $-c$. Assuming that the left door is reached
after $n$ steps and gives reward $R_{\text{left}}=1$, while the right door is reached after $m$ steps and gives reward
$R_{\text{right}}=2$, we get for $0<\gamma<1$ (discount):

\[V_{\text{left}} = -c \frac{1-\gamma^n}{1-\gamma} + \gamma^n\]
\[V_{\text{right}} = -c \frac{1-\gamma^m}{1-\gamma} + 2\gamma^m\]
When setting it up using the geometric-series formula and discounted return definition. From this we get
\[
V_{\text{left}} - V_{\text{right}}
=
-c\frac{1-\gamma^n}{1-\gamma} + \gamma^n
+ c\frac{1-\gamma^m}{1-\gamma} - 2\gamma^m
\]
\[
= \gamma^n\left(1+\frac{c}{1-\gamma}\right)
- \gamma^m\left(2+\frac{c}{1-\gamma}\right)
\]
Because
\[
V_{\text{left}} > V_{\text{right}}
\Leftrightarrow
\gamma^n\left(1+\frac{c}{1-\gamma}\right)
>
\gamma^m\left(2+\frac{c}{1-\gamma}\right),
\]
we see that if
\[
\gamma^n\left(1+\frac{c}{1-\gamma}\right)
>
\gamma^m\left(2+\frac{c}{1-\gamma}\right)
\]
best option is to go left, if
\[
\gamma^n\left(1+\frac{c}{1-\gamma}\right)
<
\gamma^m\left(2+\frac{c}{1-\gamma}\right)
\]
best option is to go right, and when
\[
\gamma^n\left(1+\frac{c}{1-\gamma}\right)
=
\gamma^m\left(2+\frac{c}{1-\gamma}\right)
\]
it makes no difference.


%Notes:
%Reward right door: 2
%Reward left door: 1

%Cost step: > 0

%n tiles on left
%m tiles on right

%gamma: discount


%$$
%V_{left} = (-cost) \times \frac{1 - \gamma^n}{1 - \gamma} + \gamma^n \times R_{left}
%$$
%$$
%V_{right} = (-cost) \times \frac{1 - \gamma^m}{1 - \gamma} + \gamma^m \times R_{right}
%$$

%It requires $V_{left} = V_{right}$ for equilibrium


\end{document}